% !TeX root = ../document_maitre.tex

\begin{itemize}
	\item Présentations des résultats obtenus pour chaque fonctionnalité/technologies implémentées. Ce n’est pas dit que je puisse avoir de résultats concrets à présenter, auquel cas je me concentrerais peut-être plus sur les résultats espérés et les premiers retours si jamais nous avons interrogé des usagers de la solution sur ce qu’ils attendent liées à ces fonctionnalités, ce qu'ils en attendent, etc. 
	% \item  Présentations des fonctionnalités créées et le parcours pensé. 
	% \item  Espoirs et réalités : quelles différences entre le point de départ de la réflexion sur ces fonctionnalités et ce qui a abouti ? Pourquoi ? 
	% \item Résultats 
	
	\item \textbf{Plan :}
	\subitem 1. La recherche augmentée 
	\subsubitem 1.1 Présentation du processus technique imaginé et/ou établi, si on en est là, pour l'aspect recherche augmentée. Explication du point de départ, de l'angle de travail suivi, et de chaque décision
	\subsubitem 1.2 Explication des défis relevés, en relation avec le Chapitre 3, et des solutions apportées aux points bloquants identifiés.
	\subsubitem 1.3 Parcours utilisateurs pensés et en quoi ils respectent des principes d'UX design et à intégrer cette pensée pour faciliter l'usage.
	\subsubitem 1.4 Retracer les différents changements importants, qu'ils soient techniques ou UX Design. Pourquoi ce changement ? Quelles solutions apportées ? Difficultés rencontrées dans l'établissement de la solution ? 
	
	\subitem 2. Le "Chatbot"
	\subsubitem 2.1 Présentation du processus technique imaginé et/ou établi, si on en est là, pour l'aspect recherche augmentée. 
	\subsubitem 2.2 Explication des défis relevés, en relation avec le Chapitre 4, et des solutions apportées aux points bloquants identifiés.
	\subsubitem 2.3 Parcours utilisateurs pensés et en quoi ils respectent des principes d'UX design et à intégrer cette pensée pour faciliter l'usage.
	\subsubitem 2.4 Retracer les différents changements importants, qu'ils soient techniques ou UX Design. Pourquoi ce changement ? Quelles solutions apportées ? Difficultés rencontrées dans l'établissement de la solution ?
	
	\subitem 3. Analyses d'écarts entre l'espoir d'origine, les besoins utilisateurs si exprimés, et la solution finale. Explications des différences et de leurs raisons d'êtres \textit{(relation avec les sections 2 et 3 du chapitre 2 de la partie 1)}
	
	\subitem 4. Retours utilisateurs et analyse des retours si possible. 
	
\end{itemize}

\subsection{La recherche augmentée :}

\subsubsection{Conception techniques finale}

\subsubsection{Conception du parcours utilisateurs final}

\subsubsection{Analyse réflexive}

\subsection{Le \gls{chatbot} :}

\subsubsection{Conception techniques finale}

\subsubsection{Conception du parcours utilisateurs final}

\subsubsection{Analyse réflexive}